% ch1_answers.tex -- Quick check answers, Module 6 Chapter 1.
\begin{chapteranswers}{1}{Charged Particles in Electric and Magnetic Fields}
  % ---- 1.1 Charges in a uniform electric field
  \item The charge is $q = 2e = 2 \times \num{1.602e-19} = \qty{3.204e-19}{C}$
        and $V = \qty{2000}{V}$. So
        $W = qV = \num{3.204e-19} \times 2000 = \qty{6.408e-16}{J}$.
        In electron volts: a charge of $2e$ through \qty{2000}{V} gains
        $2 \times 2000 = \qty{4000}{eV}$.
        Answer: \final{$W = \qty{6.4e-16}{J} = \qty{4.0e3}{eV}$}.
  \item From rest, $qV = \frac{1}{2}mv^2$. Make $V$ the subject:
        \[ V = \frac{mv^2}{2q} = \frac{\num{9.109e-31} \times (\num{3.0e6})^2}{2 \times \num{1.602e-19}}
           = \qty{25.6}{V}. \]
        Answer: \final{$V = \qty{26}{V}$}.
  \item (a) The weight acts down, so the electric force must act up. The drop
        is negative, so the electric force is opposite to the field. So the
        field points down, from the top plate to the bottom plate: the
        \textbf{top plate is positive}.
        (b) The forces balance: $qE = mg$, with $E = \dfrac{V}{d}$. So
        $q\dfrac{V}{d} = mg$ and
        \[ q = \frac{mgd}{V} = \frac{\num{2.4e-15} \times 9.8 \times \num{6.0e-3}}{294}
           = \qty{4.80e-19}{C}. \]
        Number of extra electrons: $n = \dfrac{q}{e} = \dfrac{\num{4.80e-19}}{\num{1.602e-19}}
        = 3.0$. Answer: \final{$q = \qty{4.8e-19}{C}$, 3 extra electrons}.
  \item The student is \textbf{wrong}. With the same voltage, a bigger gap $d$
        gives a weaker field, $E = \dfrac{V}{d}$, so a smaller force $qE$. The
        smaller force acts over a longer distance. The work is
        $W = qEd = q\left(\dfrac{V}{d}\right)d = qV$, which does not change. So
        the kinetic energy, and the final speed, stay the same.
  % ---- 1.2 Curved paths in an electric field
  \item The proton is positive, so it is pushed along the field: \textbf{down},
        towards the negative plate. The field is
        $E = \dfrac{V}{d} = \dfrac{600}{0.030} = \qty{2.0e4}{V.m^{-1}}$.
        Its acceleration is
        $a = \dfrac{eE}{m_\mathrm{p}} = \dfrac{\num{1.602e-19} \times \num{2.0e4}}{\num{1.673e-27}}
        = \qty{1.915e12}{m.s^{-2}}$. The time between the plates is the same as
        before, $t = \qty{2.5e-9}{s}$ (same speed, same length). So
        $y = \frac{1}{2}at^2 = \frac{1}{2} \times \num{1.915e12} \times (\num{2.5e-9})^2
        = \qty{5.98e-6}{m}$.
        Answer: \final{$y = \qty{6.0e-6}{m}$ downwards}. This is about
        $\frac{1}{1840}$ of the electron's deflection, because the proton's mass
        is about 1840 times bigger.
  \item Similarities (any two): (i) both paths are parabolas; (ii) both have a
        constant velocity in one direction and a constant acceleration at right
        angles to it; (iii) both are solved with $s = ut + \frac{1}{2}at^2$ and
        $v = u + at$ in each direction.
        Differences (any two): (i) the electron's acceleration $\dfrac{qE}{m}$
        depends on its charge and mass, whereas $g$ is the same for every mass;
        (ii) the electron is pushed towards the positive plate (the direction
        depends on the sign of the charge), whereas the ball always falls down;
        (iii) the electron's acceleration is usually far bigger than $g$ (about
        \num{1.4e15} times $g$ in \cref{ch1:we:accel}, and \num{3.6e14} times
        $g$ in \cref{ch1:we:deflect}), so its path lasts only nanoseconds; (iv) the electric force acts only on charges,
        whereas gravity acts on every mass.
  \item Doubling $u$ halves the time between the plates:
        $t = \dfrac{L}{u} = \dfrac{0.060}{\num{4.8e7}} = \qty{1.25e-9}{s}$. The
        acceleration is the same, $a = \qty{3.517e15}{m.s^{-2}}$. So
        $y = \frac{1}{2}at^2 = \frac{1}{2} \times \num{3.517e15} \times (\num{1.25e-9})^2
        = \qty{2.748e-3}{m}$. Answer: \final{$y = \qty{2.7}{mm}$}. This is a
        quarter of the old deflection
        ($\qty{1.099}{cm} \div 4 = \qty{0.2748}{cm}$), because $y \propto t^2$
        and $t$ halved.
  \item Take up as positive. The force is down, so the acceleration is
        $a = -\dfrac{qE}{m} = -\dfrac{qV}{md}$. The starting vertical velocity
        is $u_y = u\sin\theta$. When the particle lands, its vertical
        displacement is zero:
        \begin{align*}
          0 &= (u\sin\theta)\,t - \frac{1}{2}\,\frac{qV}{md}\,t^2
            && \text{$s = ut + \tfrac{1}{2}at^2$ with $s = 0$}\\
          0 &= u\sin\theta - \frac{1}{2}\,\frac{qV}{md}\,t
            && \text{divide by $t$ ($t \ne 0$)}\\
          \frac{qV}{2md}\,t &= u\sin\theta
            && \text{move the $t$ term to the left}\\
          t &= \frac{2mdu\sin\theta}{qV}
            && \text{multiply both sides by $\tfrac{2md}{qV}$}
        \end{align*}
        This is the same as $t = \dfrac{2u_y}{a}$ in \cref{ch1:we:angled}, with
        $a = \dfrac{qV}{md}$ (the size of the acceleration).
  % ---- 1.3 Charges in a uniform magnetic field
  \item For a positive charge: fingers out of the page (field), thumb towards
        the top of the page (velocity); the palm faces to the right. An electron
        is negative, so reverse this. Answer: \final{towards the left}.
  \item The velocity is in the plane of the page and the field is into the
        page. So the angle between $v$ and $B$ is \ang{90}, not \ang{35}:
        $F = qvB = \num{1.602e-19} \times \num{5.0e5} \times 0.40 = \qty{3.2e-14}{N}$.
        Direction: fingers into the page, thumb along $v$; the palm faces
        \ang{90} anticlockwise from $v$. Answer: \final{$F = \qty{3.2e-14}{N}$}.
        The force is in the plane of the page, at right angles to $v$. It
        points up and to the left. Its angle is \ang{125}, measured
        anticlockwise from the right. (Using $\sin\ang{35}$ would wrongly give
        \qty{1.8e-14}{N}.)
  \item (a) Set $qvB = \dfrac{mv^2}{r}$, so
        $B = \dfrac{mv}{qr} = \dfrac{\num{1.673e-27} \times \num{2.0e6}}{\num{1.602e-19} \times 0.50}
        = \qty{4.18e-2}{T}$. Answer: \final{$B = \qty{4.2e-2}{T}$} (\qty{42}{mT}).
        (b) $r = \dfrac{mv}{qB}$, so $r \propto v$: the \textbf{radius doubles}.
        $T = \dfrac{2\pi m}{qB}$ has no $v$ in it: the \textbf{period stays the
        same}. The circle's circumference doubles, but the proton also moves
        twice as fast, so one turn still takes the same time:
        $T = \dfrac{2\pi r}{v} = \dfrac{2\pi \times 0.50}{\num{2.0e6}}
        = \dfrac{2\pi \times 1.0}{\num{4.0e6}} = \qty{1.6e-6}{s}$ both times.
  \item Split the velocity into two parts. The part across the field,
        $v\sin\ang{40}$ (about 64\,\% of $v$), feels a force at right angles to
        it, so it gives circular motion. The part along the field,
        $v\cos\ang{40}$ (about 77\,\% of $v$), feels no force, so it stays
        constant and carries the electron along the field lines. Together the
        path is a \textbf{helix} (a spiral) along the field. The speed stays
        constant, because the force does no work.
  % ---- 1.3 (subsection) Comparing electric, magnetic and gravitational fields
  \item Both move in circles at constant speed. Their charges are the same size,
        so $r = \dfrac{mv}{qB} \propto m$: the proton's circle is about 1840 times
        bigger ($\dfrac{m_\mathrm{p}}{m_\mathrm{e}} = 1837$). Their charges have
        opposite signs, so the forces are opposite and they turn opposite ways
        (with $B$ into the page: proton anticlockwise, electron clockwise).
        Only an \textbf{electric} field could change their kinetic energy. The
        electric force can act along the motion, so it does work ($W = qV$). The
        magnetic force is always perpendicular to the velocity, so
        $W = Fs\cos\ang{90} = 0$.
  \item Similarity: both are uniform circular motion. In both, the force
        towards the centre is perpendicular to the velocity. So it does no work,
        and the speed stays constant. Differences (any two): (i) the satellite's
        centripetal force is gravity, $F = \dfrac{GMm}{r^2}$, whereas the
        proton's is the magnetic force $qvB$, which exists only because the
        proton moves; (ii) gravity depends on the distance from the Earth,
        whereas $qvB$ depends on the speed; (iii) a bigger orbit is slower and
        has a longer period, whereas for the proton $r \propto v$ and the
        period $T = \dfrac{2\pi m}{qB}$ does not depend on the speed.
  \item (a) Electric: $F = qE = \num{1.602e-19} \times \num{1.0e4}
        = \qty{1.602e-15}{N}$. Magnetic ($\theta = \ang{90}$, so
        $\sin\theta = 1$): $F = qvB = \num{1.602e-19} \times \num{2.0e5}
        \times 0.050 = \qty{1.602e-15}{N}$.
        Answer: \final{both forces are \qty{1.6e-15}{N}}. They are equal only
        at this speed. At twice the speed the magnetic force doubles, but the
        electric force does not change.
        (b)~Only the \textbf{electric} force. $F = qE$ does not depend on the
        speed. The magnetic force $qvB$ is zero when $v = 0$.
\end{chapteranswers}
